\documentclass[11pt,a4paper]{article}

\usepackage[utf8]{inputenc}
\usepackage[T1]{fontenc}
\usepackage{lmodern}
\usepackage[polish]{babel}
\usepackage{polski}
\usepackage{geometry}
\geometry{
    top=2.5cm,
    bottom=2.5cm,
    left=2.5cm,
    right=2.5cm
}

\usepackage{amsmath,amssymb,amsfonts}
\usepackage{graphicx}
\usepackage{booktabs}
\usepackage{subcaption}
\usepackage{float}
\usepackage{microtype}
\usepackage[hidelinks]{hyperref}

\begin{document}

\pagenumbering{gobble}

% ==============================================================================
% STRONA TYTUŁOWA
% ==============================================================================
\begin{titlepage}
    \centering
    \vspace*{1.5cm}
    
    {\scshape\Large Politechnika Warszawska \par}
    \vspace{0.3cm}
    {\scshape\large Wydział Matematyki i Nauk Informacyjnych \par}
    
    \vspace{2.5cm}
    
    {\huge\bfseries Projekt 1 - część druga\\[0.3em]Klasyfikacja zbioru MNIST wielowarstwowym perceptronem\par}
    
    
    \vspace{0.8cm}
    {\Large\itshape Przedmiot: Sieci Neuronowe \par}
    
    \vspace{3.0cm}
    
    \begin{minipage}[t]{0.48\textwidth}
        \flushleft
        \textbf{Autorzy:}\\[0.3em]
        Jakub Rak\\
        Krzysztof Radzewicz
    \end{minipage}
    \hfill
    \begin{minipage}[t]{0.48\textwidth}
        \flushright
        \textbf{Prowadzący:}\\[0.3em]
        mgr inż. Mikołaj Małkiński
    \end{minipage}
    \vfill
    
    {\large Warszawa, \today \par}
\end{titlepage}

\pagenumbering{roman}
\tableofcontents
\newpage

\pagenumbering{arabic}

\section{Opis zbioru danych MNIST}

Niniejszy raport stanowi uzupełnienie pierwszego projektu laboratoryjnego.
W pierwszej części model \textbf{MLP} testowany był na syntetycznych zbiorach dwuwymiarowych
(klasyfikacja $\mathbb{R}^2 \to \{0, 1\}$ oraz $\mathbb{R}^2 \to \{0, 1, 2\}$) oraz regresji jednowymiarowej ($\mathbb{R} \to \mathbb{R}$).
W tej części oceniono możliwość skalowania zaimplementowanej sieci neuronowej do danych o większym wymiarze wejściowym -
klasycznym zbiorze \textit{MNIST} w zadaniu klasyfikajci odręcznie pisanych cyfr.

Zbiór \textit{MNIST} zawiera łącznie 70~000 znormalizowanych, czarno-białych (zapisanych w 256 poziomach szarości) obrazów przedstawiających cyfry arabskie od 0 do 9.
Dane rozbito na dwa zestawy:
\begin{itemize}
    \item \textbf{Oficjalny zbiór treningowy}: 60~000 obrazów, z których w eksperymencie 55~000 służyło do optymalizacji parametrów, a 5~000 do walidacji,
    \item \textbf{Zbiór testowy}: 10~000 niezależnych obrazów służących do ostatecznej ewaluacji sieci.
\end{itemize}

Każda próbka to obrazek o stałych wymiarach $28 \times 28$ pikseli.
Każdy piksel reprezentowany jest przez liczbę całkowitą z zakresu $[0, 255]$, gdzie wartość 0 oznacza tło,
a 255 maksymalną intensywność piksela.

W eksperymencie obrazy poddano spłaszczeniu do jednowymiarowego wektora cech:
\begin{equation}
    \mathbf{x} \in \mathbb{R}^{784}, \quad D = 28 \times 28 = 784.
\end{equation}
Wartości pikseli zostały znormalizowane do przedziału $[0, 1]$.
Etykiety klas reprezentują 10 wzajemnie wykluczających się kategorii
($y \in \{0, 1, \dots, 9\}$), które zakodowano w postaci wektorów \textit{one-hot} o długości 10.

\section{Zmiany w kodzie}

Implementacja modelu \textbf{MLP} nie uległa żadnym modyfikacjom. Dodano jedynie dwa moduły pomocnicze:
\begin{itemize}
    \item \textbf{Wczytywanie danych} (\texttt{experiments/common/data\_loader.py}): funkcja odczytująca binarne pliki IDX zbioru MNIST,
    spłaszczająca obrazy $28 \times 28$ do wektorów o długości 784, normalizująca wartości pikseli do przedziału $[0, 1]$ oraz kodująca etykiety w formacie \textit{one-hot}.
    \item \textbf{Wizualizacja wyników} (\texttt{experiments/common/plotting.py}): funkcje do wykreślania krzywych uczenia oraz generowania macierzy pomyłek wraz z metrykami ewaluacyjnymi na zbiorze testowym.
\end{itemize}

\section{Przebieg eksperymentu i analiza wyników}

\subsection{Konfiguracja modelu i procedura uczenia}

Eksperyment zrealizowano w skrypcie \texttt{experiments/run\_mnist.py}.
Krzywe dynamiki uczenia uśredniono po 10 uruchomieniach z ziarnami losowości
\texttt{54217--54226}. Ziarna te sterowały inicjalizacją wag i tasowaniem mini-batchy.
Podział na próbki uczące i walidacyjne był taki sam dla wszystkich uruchomień
(ziarno podziału: \texttt{54217}), dzięki czemu pasma zmienności opisują wpływ losowości treningu,
a nie zmiany składu podzbiorów. Tabela metryk, macierz pomyłek i przykłady błędów
w dalszej części raportu dotyczą pojedynczego modelu z dotychczasowym ziarnem \texttt{54217}.
Z oficjalnego zbioru treningowego (60~000 próbek) losowo wydzielono 5~000 próbek na podzbiór walidacyjny (służył do wyświetlania
realnych postępów w uczeniu), a pozostałe 55~000 próbek przeznaczono do treningu. Niezależny zbiór testowy liczył 10~000 próbek.
Liczebność testu potwierdzono w nagłówkach obu lokalnych plików IDX
(\texttt{t10k-images.idx3-ubyte} i \texttt{t10k-labels.idx1-ubyte});
ewaluacja korzystała z całego tego zbioru, bez opcji ograniczającej liczbę próbek.

Wykorzystano następującą architekturę: $[784, 150, 50, 10]$.
\begin{itemize}
    \item \textbf{Aktywacje ukryte}: wszystkie warstwy ukryte z aktywacją \texttt{ReLU} i inicjalizacją wag \textit{He},
    \item \textbf{Warstwa wyjściowa}: 10 neuronów z aktywacją \texttt{Softmax}, odpowiadających prawdopodobieństwom przynależności do klas $0, 1, \dots, 9$.
\end{itemize}
We wszystkich warstwach włączono wektory \textit{bias}.
Jako funkcję kosztu zastosowano kategoryczną entropię krzyżową (\texttt{CategoricalCrossEntropyLoss}).
Optymalizację prowadzono za pomocą klasycznego \texttt{SGD} ze stałym współczynnikiem uczenia $\eta = 0.05$.
Sieć trenowano przez 30 epok z rozmiarem mini-batcha równym 32.

\subsection{Dynamika uczenia — krzywe straty i dokładności}

Na Rysunku~\ref{fig:mnist_learning_curves} przedstawiono przebieg uczenia modelu na przestrzeni 30 epok.
Każda linia przedstawia średnią danej metryki w danej epoce z 10 ziaren,
a półprzezroczyste pasmo zakres średnia $\pm$ jedno odchylenie standardowe.
Odchylenie obliczono jako $\sigma_e = \sqrt{\frac{1}{10}\sum_{s=1}^{10}(m_{s,e}-\bar{m}_e)^2}$,
gdzie $m_{s,e}$ jest wartością metryki dla ziarna $s$ w epoce $e$.
Pasmo opisuje rozrzut wyników między uruchomieniami, nie przedział ufności średniej.

\begin{figure}[H]
    \centering
    \includegraphics[width=0.98\textwidth]{../experiments/results/mnist/mnist_learning_curves_mean_std.png}
    \caption{Przebieg uczenia modelu MLP $[784, 150, 50, 10]$ na zbiorze MNIST: strata kategorycznej entropii krzyżowej (panel lewy) oraz dokładność klasyfikacji w procentach (panel prawy) na zbiorze uczącym i walidacyjnym. Linie: średnie z 10 ziaren; pasma: $\pm 1\sigma$ między ziarnami.}
    \label{fig:mnist_learning_curves}
    \end{figure}

Analiza krzywych uczenia pozwala sformułować następujące obserwacje:
\begin{itemize}
    \item \textbf{Funkcja straty}: W lewym panelu widać gwałtowny spadek straty w pierwszych 5 epokach.
    Średnia strata treningowa zbliża się do zera. Stratę walidacyjną należy interpretować razem z pasmem zmienności:
    jej wzrost mimo dalszego spadku straty treningowej jest oznaką przeuczenia według tej metryki,
    nawet jeśli dokładność walidacyjna pozostaje wysoka.
    Minimum średniej straty walidacyjnej przypada na 17. epokę; w 30. epoce wynosi ona
    $0.0782 \pm 0.0045$, przy stracie treningowej około $0.00124$.
    \item \textbf{Krzywa dokładności}: Dokładność rośnie szybko już w pierwszych epokach. Jest tak ponieważ wagi są aktualizowane
    co mini batch, co w połączeniu z bardzo dużym zbiorem treningowym (55~000 próbek) i rozmiarem mini batch 32, daje bardzo wiele aktualizacji.
    Pod koniec treningu dokładność ucząca jest bliska $100\%$.
    W 30. epoce średnia dokładność walidacyjna wynosi $97.976\%$,
    a jej odchylenie standardowe między ziarnami to $0.150$ punktu procentowego.
    \item \textbf{Zmienność między ziarnami}: Pasma pokazują wrażliwość procesu uczenia na inicjalizację wag
    i kolejność mini-batchy przy stałym podziale danych. Pozwala to ocenić powtarzalność przebiegu uczenia
    bez opierania wniosków wyłącznie na pojedynczym uruchomieniu.
\end{itemize}

\subsection{Ewaluacja na zbiorze testowym i macierz pomyłek}

Ostateczna ewaluacja modelu dla ziarna \texttt{54217} na 10~000 niezależnych próbkach testowych wykazała całkowitą dokładność na poziomie \textbf{98.00\%}
(9~800 poprawnie sklasyfikowanych obrazów na 10~000).
Szczegółowe metryki w rozbiciu na poszczególne cyfry zebrano w Tabeli~\ref{tab:mnist_metrics}.

\begin{table}[H]
    \centering
    \caption{Szczegółowe metryki jakości klasyfikacji modelu MLP na zbiorze testowym MNIST ($N = 10~000$).}
    \label{tab:mnist_metrics}
    \resizebox{\textwidth}{!}{%
    \begin{tabular}{crrrrrr}
        \toprule
        \textbf{Klasa} & \textbf{Liczba wystąpień} & \textbf{Recall (TPR)} & \textbf{Precision (PPV)} & \textbf{Specificity (TNR)} & \textbf{FPR} & \textbf{$F_1$-Score} \\
        \midrule
        0 & 980  & 0.9898 & 0.9838 & 0.9982 & 0.0018 & 0.9868 \\
        1 & 1135 & 0.9912 & 0.9912 & 0.9989 & 0.0011 & 0.9912 \\
        2 & 1032 & 0.9777 & 0.9825 & 0.9980 & 0.0020 & 0.9801 \\
        3 & 1010 & 0.9822 & 0.9706 & 0.9967 & 0.0033 & 0.9764 \\
        4 & 982  & 0.9786 & 0.9776 & 0.9976 & 0.0024 & 0.9781 \\
        5 & 892  & 0.9742 & 0.9797 & 0.9980 & 0.0020 & 0.9770 \\
        6 & 958  & 0.9812 & 0.9822 & 0.9981 & 0.0019 & 0.9817 \\
        7 & 1028 & 0.9767 & 0.9814 & 0.9979 & 0.0021 & 0.9790 \\
        8 & 974  & 0.9764 & 0.9814 & 0.9980 & 0.0020 & 0.9789 \\
        9 & 1009 & 0.9703 & 0.9683 & 0.9964 & 0.0036 & 0.9693 \\
        \midrule
        \multicolumn{2}{l}{\textbf{Średnia arytmetyczna (makro)}} & 0.9798 & 0.9799 & 0.9978 & 0.0022 & 0.9798 \\
        \multicolumn{2}{l}{\textbf{Średnia ważona ($N = 10~000$)}} & 0.9800 & 0.9800 & 0.9978 & 0.0022 & 0.9800 \\
        \bottomrule
    \end{tabular}%
    }
\end{table}

Strukturę błędów popełnianych przez sieć szczegółowo ilustruje macierz pomyłek zaprezentowana na Rysunku~\ref{fig:mnist_confusion_matrix}.

\begin{figure}[H]
    \centering
    \includegraphics[width=0.72\textwidth]{../experiments/results/mnist/mnist_confusion_matrix.png}
    \caption{Znormalizowana macierz pomyłek na zbiorze testowym MNIST ($N = 10~000$). Wiersze odpowiadają etykietom prawdziwym, natomiast kolumny predykcjom modelu.}
    \label{fig:mnist_confusion_matrix}
    \end{figure}

Z analizy macierzy pomyłek i tabeli metryk płyną konkretne wnioski:
\begin{itemize}
    \item \textbf{Najłatwiejsze klasy}: Model najskuteczniej rozpoznaje cyfrę \textit{1}
    ($F_1 = 0.9912$, 1125 trafień na 1135 próbek) oraz cyfrę \textit{0} ($F_1 = 0.9868$, 970 trafień na 980 próbek).
    Obie te cyfry są bardzo proste geometrycznie, co prawdopodobnie redukuje liczbę pomyłek.
    \item \textbf{Przykłady błędnie sklasyfikowanych próbek}: Rysunek~\ref{fig:mnist_misclassified} przedstawia trzy przykładowe próbki ze zbioru testowego,
    dla których model podjął błędną decyzję: cyfrę 4 sklasyfikowaną jako 9, cyfrę 5 jako 3 oraz cyfrę 7 jako 2.
    
    \begin{figure}[H]
        \centering
        \includegraphics[width=0.88\textwidth]{../experiments/results/mnist/mnist_misclassified_examples.png}
        \caption{Przykładowe obrazy ze zbioru testowego MNIST błędnie sklasyfikowane przez model MLP, obrazujące główne pomyłki modelu.}
        \label{fig:mnist_misclassified}
    \end{figure}
    
    \item \textbf{Bardzo niska liczba fałszywych alarmów}: Wartości \textit{TNR} dla wszystkich klas wynoszą
    co najmniej $0.9964$, a odsetek fałszywych alarmów \textit{FPR} nie przekracza $0.36\%$,
    co wskazuje na niezawodność aktywacji Softmax bez faworyzowania żadnej z kategorii.
\end{itemize}

\section{Krótkie podsumowanie i wnioski}

Przeprowadzony eksperyment na zbiorze MNIST pozwolił zweryfikować zachowanie zaimplementowanego
wielowarstwowego perceptronu w znacznie bardziej wymagających warunkach  niż proste dane dwuwymiarowe.
Główne wnioski można podsumować następująco:

\begin{enumerate}
    \item \textbf{Poprawność i skalowalność autorskiej implementacji}:
    Biblioteka \texttt{py\_nn1} bez żadnych modyfikacji
    poradziła sobie z optymalizacją modelu posiadającego ponad 125~000 parametrów wagowych
    ($784 \times 150 + 150 \times 50 + 50 \times 10 \approx 125~850$ wag wraz z biasami)
    na zbiorze 55~000 próbek uczących. Obliczenia przebiegły stabilnie, bez problemów
    zanikających czy eksplodujących gradientów.
    
    \item \textbf{Wysoka skuteczność klasycznego MLP}:
    Standardowy perceptron z dwiema warstwami ukrytymi ($150$ i $50$ neuronów),
    nieliniowością \texttt{ReLU} i inicjalizacją \textit{He} uzyskał \textbf{98.00\% dokładności}
    na 10~000 obrazów testowych.
    
    \item \textbf{Sprawność optymalizatora SGD}:
    Prosty algorytm stochastycznego spadku gradientu z $\eta = 0.05$
    i wielkością mini batcha 32 okazał się w pełni wystarczający do osiągnięcia zbieżności
    w granicach 15--20 epok, zapewniając jednocześnie odpowiednią jakość predykcji.
    
    \item \textbf{Błędy klasyfikacji}:
    Można przypuszczać że spośród nielicznych 200 błędów (na 10~000 przypadków testowych)
    istotna część z nich wynika z podobieństwa geometrycznego między klasami cyfr.
    Warto byłoby zbadać czy takie błędy można ograniczyć
    stosując model który zamiast spłaszczonego wektora, przyjmowałby wejście zachowujące dwuwymiarowe zależności
    (na przykład konwolucyjną sieć neuronową).
\end{enumerate}

\end{document}
